\subsection{Prime Density}

\subsubsection{Prime Counting Function}

The prime counting function $\pi(n)$ denotes the number of prime numbers less than or equal to $n$.

\subsubsection{Prime Number Theorem}

The Prime Number Theorem provides an asymptotic approximation for $\pi(n)$:

\begin{equation}
    \pi(n) \sim \frac{n}{\ln(n)}
\end{equation}

More precisely:
\begin{equation}
    \lim_{n \to \infty} \frac{\pi(n)}{n/\ln(n)} = 1
\end{equation}

\subsubsection{Better Approximations}

A more accurate approximation is given by the logarithmic integral:

\begin{equation}
    \pi(n) \sim \text{Li}(n) = \int_2^n \frac{dt}{\ln(t)}
\end{equation}

For practical purposes, the following approximation is also useful:

\begin{equation}
    \pi(n) \approx \frac{n}{\ln(n) - 1} \quad \text{for } n \geq 17
\end{equation}

\subsubsection{Asymptotic Density}

The probability that a random integer near $n$ is prime is approximately $\frac{1}{\ln(n)}$.

The $n$-th prime number $p_n$ can be approximated by:

\begin{equation}
    p_n \sim n \ln(n)
\end{equation}
